\begin{abstract}
Honesty is a fundamental principle for aligning large language models (LLMs) with human values, requiring these models to recognize what they know and don't know and be able to faithfully express their knowledge. Despite promising, current LLMs still exhibit significant dishonest behaviors, such as confidently presenting wrong answers or failing to express what they know. In addition, advancing research on the honesty of LLMs requires addressing challenges, including varying definitions of honesty, difficulties in distinguishing between known and unknown knowledge, and a lack of comprehensive understanding of related research. To address these issues, we provide a survey on the honesty of LLMs, covering its clarification, evaluation approaches, and strategies for improvement. Moreover, we offer insights for future research, aiming to inspire further exploration in this important area.
\end{abstract}
